\documentclass[11pt]{article}
\usepackage[margin=1.15in]{geometry}
\usepackage{amsmath,amssymb,amsthm}
\usepackage{hyperref}
\sloppy
\usepackage{xcolor}
\newtheorem{theorem}{Theorem}
\newtheorem{proposition}[theorem]{Proposition}
\newtheorem{lemma}[theorem]{Lemma}
\newtheorem{corollary}[theorem]{Corollary}
\newtheorem{conjecture}[theorem]{Conjecture}
\theoremstyle{remark}\newtheorem{remark}[theorem]{Remark}
\DeclareMathOperator{\perm}{per}
\DeclareMathOperator{\lcm}{lcm}
\newcommand{\Z}{\mathbb Z}
\newcommand{\N}{\mathbb N}
\newcommand{\vtwo}{v_2}

\title{Machine-Verified Experimental Mathematics:\\
Four Ramanujan-Style Case Studies}
\author{Vico Bonfioli\thanks{All data, scripts, and the
Lean~4 development are in the \texttt{rama-notes} repository; the formalization is the
library \texttt{RamaLean/} (Lean~4 + Mathlib, standard axioms
\texttt{[propext, Classical.choice, Quot.sound]}; plus \texttt{Lean.ofReduceBool}
for the \texttt{native\_decide} computational cross-checks (Cases~1--2), no
\texttt{sorry}. The Paper~3 divisibility theorems and the permanent library are on
the clean three-axiom base). Correspondence: \texttt{vicobonfioli@gmail.com}.}}
\date{July 2026}

\begin{document}
\maketitle

\begin{abstract}
Ramanujan's method---compute aggressively, conjecture boldly from the data, then
seek the identity---produced curiosities that were later revealed as tips of deep
icebergs. We ask a \emph{methodological} question: can that method be run today
end-to-end with machine-checked rigor, so that ``conjecture from computation''
terminates in ``theorem verified by kernel''? We answer yes through four
self-contained case studies: three of the four reach a Lean~4/Mathlib proof, and
Case~4 is computation plus paper-proof, resting on the shared formalized
matching-polynomial identities. The contribution is \emph{not} any individual theorem---we
argue plainly that the four discoveries are, in the grand scheme, minor---but the
demonstration that the experimental-to-formal pipeline closes, including (to our
knowledge) the first machine-verified instance of the Hall--Puder--Sawin
``Ramanujan coverings'' circle of ideas, and a small library of reusable formal
permanent theory that Mathlib previously lacked. We are equally explicit about the
pipeline's failure modes: one case study is presented as a cautionary tale in
which an intermediate ``decisive'' computation was \emph{retracted} after its
proxy invariant was found unfaithful, leaving a razor-sharp but still-open
conjecture together with the single experiment that would settle it.
\end{abstract}

\section{Introduction}

Ramanujan filled notebooks with identities found by computing far past the point
where a pattern becomes ``obvious,'' then trusting the pattern. A century of
hindsight shows how often the obvious pattern hid real structure---partitions and
modular forms, the mock theta functions and mock modular forms. The working method
was: \emph{compute, conjecture, and only then prove.}

Two things have changed since. First, ``compute'' can now mean tens of billions of
exact operations, pushing the empirical frontier of a combinatorial sequence far
enough that the eventual asymptotic regime is actually visible. Second,
``prove'' can now mean \emph{machine-checked}: a proof assistant's kernel certifies
the argument against a tiny trusted core. This paper is an experiment in running
Ramanujan's loop with both upgrades in place, on four genuinely new (if modest)
objects, and reporting honestly on what the discipline of formalization does and
does not buy.

Our thesis is deliberately narrow. We do \emph{not} claim the four objects below
are deep; Section~\ref{sec:assessment} argues the opposite. We claim that the
\emph{pipeline}
\[
   \text{compute} \;\to\; \text{conjecture} \;\to\; \text{structural proof}
   \;\to\; \text{Lean formalization} \;\to\; \text{decisive experiment when theory
   stalls}
\]
closes end-to-end on real examples, that closing it surfaces reusable
infrastructure, and that it imposes a useful honesty: a formalized theorem cannot
be quietly wrong, and---as Case~3 shows---a computation engineered to
settle a conjecture can force a \emph{retraction} rather than a confirmation.

\section{The pipeline}\label{sec:pipeline}

Two features distinguish the loop from ordinary formalization. (i) The conjectures
are \emph{discovered} here, not imported from the literature to be re-proved.
(ii) When a structural proof stalls, we engineer a \emph{decisive computation}
pushed to the feasibility frontier (Gray-code Ryser for permanents; sieve methods
to $10^7$ for partitions), and we allow that computation to \emph{refute}. The
formalization step is not decoration: it is where an argument that ``reads
correctly'' is forced to be correct, and several of our human-written first drafts
did not survive it.

\section{Case study 2: random lifts of cycles, machine-verified}\label{sec:cycle}

Let $C_n$ be the $n$-cycle and consider a uniformly random permutation
$r$-\emph{lift} (an $r$-sheeted covering given by assigning a permutation in $S_r$
to each edge). The expected characteristic polynomial of the ``new'' part of the
lift---the part orthogonal to the base spectrum---is, by the
Hall--Puder--Sawin theory of Ramanujan coverings~\cite{HPS},
\begin{equation}\label{eq:cycle}
   \Phi_{C_n,r}(x) \;=\; \chi_{C_n}(x)\cdot U_{r-1}\!\big(T_n(x/2)\big),
\end{equation}
where $T_n,U_{r-1}$ are Chebyshev polynomials of the first and second kind. The
quotient $U_{r-1}(T_n(x/2))$ is the $(r-1)$-matching polynomial of $C_n$; its
real-rootedness is the finite-$r$ generalization, due to~\cite{HPS}, of the
Heilmann--Lieb theorem and of the interlacing method of Marcus, Spielman, and
Srivastava~\cite{MSS}. This structure was established in 2018; we do not improve
it. What we contribute is threefold.

\emph{An independent elementary derivation.} We obtain~\eqref{eq:cycle} through a
holonomy/transfer-matrix computation and the exponential formula, rather than
through matching-polynomial combinatorics, yielding en route the generating
function $\sum_{r\ge0}\Phi_{C_n,r}(x)\,z^r = (1-z)^2/(1-2Yz+z^2)$ with
$Y=T_n(x/2)$.

\emph{An exact, effective spectral edge.} Whereas~\cite{HPS} give the asymptotic
Ramanujan bound (largest new eigenvalue $\to 2$), we determine the edge exactly at
every finite $r$: the largest root of $U_{r-1}(T_n(x/2))$ is
\[
   x^\star \;=\; 2\cos\!\Big(\frac{\pi}{nr}\Big),
   \qquad\text{so the spectral gap is}\quad
   2 - 2\cos\!\Big(\frac{\pi}{nr}\Big) \sim \frac{\pi^2}{(nr)^2}.
\]

\emph{Full formalization.} Both the factored form~\eqref{eq:cycle} (over a general
commutative ring, with \texttt{native\_decide} cross-checks) and the spectral edge
are machine-checked in \texttt{RamaLean/Paper2*.lean}. The edge is the theorem
\begin{quote}
\texttt{Paper2SpectralEdge.spectral\_edge\_is\_root}:\quad
$U_{r-1}\big(T_n(\cos(\pi/(nr)))\big)=0$ \; for $0<n$, $2\le r$,
\end{quote}
proved from $T_n(\cos\theta)=\cos n\theta$ and the vanishing of $U_{r-1}$ at
$\cos(\pi/r)$. To our knowledge this is the first formally verified statement
inside the Ramanujan-coverings circle.

\section{Case study 3: the 2-adic gcd permanent (a cautionary tale)}\label{sec:gcd}

Let $a(n)=\perm[\gcd(i,j)]_{1\le i,j\le n}$. The determinant analogue is Smith's
1876 identity $\det[\gcd(i,j)]=\prod_{k\le n}\varphi(k)$; the permanent is catalogued
as OEIS~A085244 (2003), but its arithmetic (divisibility, 2-adic growth) appears not
to have been studied. We establish (all machine-checked): $2,3,4,8\mid a(n)$ in
suitable ranges, and the sharp linear 2-adic growth rate
\begin{equation}\label{eq:c1}
   \vtwo\big(a(n)\big) \;\ge\; n - 2\lfloor\log_2 n\rfloor - 2
   \qquad(\text{``growth constant } c=1\text{''}),
\end{equation}
the theorem \texttt{Paper3C1.two\_pow\_dvd\_permanent\_c1}. The proof rests on four
new, reusable permanent-divisibility lemmas (Section~\ref{sec:mathlib}).

\emph{The discovery, and the retraction.} Exact Gray-code Ryser values of
$\vtwo(a(n))$ at the ``peaks'' $n=2^k+1$ yield $\vtwo(a) = 3,5,11,25$ for
$k=2,3,4,5$, i.e.
\begin{equation}\label{eq:defconj}
   \vtwo\big(a(2^k+1)\big) \;=\; 2^k - 2k + 3
   \qquad\Longleftrightarrow\qquad
   D(2^k+1) := \vtwo(n!) - \vtwo(a) = 2k-4 \sim 2\log_2 n,
\end{equation}
a clean line through all four exact points, suggesting the deficit $D$ is
\emph{unbounded}, of order $\Theta(\log n)$, against the $O(\log n)$ ceiling
implied by~\eqref{eq:c1}. Full Ryser at the next peak $n=65$ is $2^{65}$
operations---out of reach---so we engineered a factored proxy
$N_0=\sum_v \perm(A_v)\perm(B_v)$ (the weight-$0$ grade of the 2-adic
decomposition $a=\sum_w 2^w N_w$), computable at $n=65$, and obtained
$\vtwo(N_0(65))=58$.

We initially read this as \emph{refuting}~\eqref{eq:defconj} (it predicts
$\vtwo(a(65))=55$). A validation check then showed the proxy is unfaithful:
$\vtwo(a)-\vtwo(N_0)$ is not identically $0$ but $-2,+1,+3$ at $n=7,12,13$. Hence
$\vtwo(N_0(65))$ does \emph{not} determine $\vtwo(a(65))$, and the ``refutation''
was withdrawn. What remains is a genuine, sharply-posed tension:
\begin{itemize}
\item the closed form~\eqref{eq:defconj} fits $4/4$ exact peaks (including $k=2$,
      where the proxy mechanism below fails), and predicts $D(65)=8$ (unbounded);
\item at $k=3,4,5$ the weight-$0$ grade dominates ($\vtwo(a)=\vtwo(N_0)$), but the
      grade-profile margins run $0,1,3,1$ (over $k=2,3,4,5$)---fluctuating, not
      growing, so this dominance is fragile; \emph{if} it nonetheless persisted,
      $\vtwo(N_0(65))=58$ would force $D(65)=5$ (a bounded ``turnover'').
\end{itemize}
These are incompatible at $k=6$: exactly one pattern breaks there. Attacking the
equivalent recursion $\vtwo(a(2^k+1))=\vtwo(a(2^{k-1}+1))+(2^{k-1}-2)$ shows it is
\emph{not} an independent, easier target---it holds iff a low grade overtakes
$w=0$ at $k=6$, iff the deficit is unbounded. What \emph{is} unconditional is the
upper bound $D(2^k+1)\le 2k$ ($k\ge3$), from~\eqref{eq:c1}; a matching growing
lower bound---equivalently the exact conjecture~\eqref{eq:defconj}---is open. A
formalized closed-form fact about the odd-$\varphi$ sub-sum
($\vtwo(A(2^k+1))=\vtwo(n!)-1$, \texttt{Paper3Atom1}) gives $D\ge2$ infinitely
often \emph{conditional} on the empirically-clear but unproven
$\vtwo(B)<\vtwo(A)$. The single decisive feasible experiment is the exact
$n=33$ grade profile (a Gray-code Ryser over $(\Z/2^{64})[t]$,
\texttt{code/rust\_grade/}); its margin between the $w=0$ grade and its nearest
competitor discriminates the two patterns.

We keep this case in the paper precisely because it is where the method behaved
best as \emph{method}: it produced a falsifiable conjecture, an intermediate
result, the discovery that the intermediate result's proxy was unfaithful, a clean
retraction, a proof that the natural next target is equivalent to the original, and
an explicit pointer to the one computation that would decide it. The formalized
core~\eqref{eq:c1} is untouched by any of this.

\section{Case studies 1 and 4 (brief)}

\emph{Case 1: partition self-divisibility.} We study $n\mid p(n)$-type divisibility
of the partition function, computed to $10^7$. A probabilistic
$\approx\sum 1/n$ heuristic (divergent) predicts infinitely many hits; one
supporting proposition is formalized. This is the shallowest case: the varying
modulus breaks the modular-forms machinery that gives Ramanujan's classical
congruences their depth, so it stays an honest curiosity.

\emph{Case 4: $d$-matching polynomial coefficients.} Inside the same HPS object of
Case~2, the coefficients $c_k(d)$ of the $d$-matching polynomial are
polynomials in graph invariants; the two leading coefficients are \emph{universal}
(functions of $|E|$ and the path count $P$ only), with $c_4$ computed exactly on
eight graphs. A modest structural refinement of a known object.

\section{The formalized library}\label{sec:mathlib}

Mathlib had no permanent-divisibility results and no permanent Laplace expansion.
Closing the pipeline produced the following, each a candidate for upstreaming:
\begin{enumerate}
\item \texttt{permanent\_succ\_column\_zero}: Laplace/cofactor expansion of the
      permanent (Mathlib had this only for \texttt{det}; the sign-free proof is
      shorter).
\item \texttt{two\_pow\_dvd\_permanent\_odd}: an all-odd matrix has
      $\vtwo(\perm)\ge \text{size}-\log$.
\item \texttt{two\_factorial\_dvd\_permanent}: two identical row-groups give
      $|A|!\,|B|! \mid \perm$, via a free $\mathrm{Sym}(A)\times\mathrm{Sym}(B)$
      action.
\item \texttt{two\_pow\_dvd\_permanent\_zeroed}: a zeroed-corner odd matrix has
      $\vtwo(\perm)\ge \text{size}-2\log$.
\item \texttt{card\_dvd\_sum\_of\_free\_invariant}: a free finite-group action with
      an invariant summand gives $|G|\mid\sum$.
\item \texttt{spectral\_edge\_is\_root} (Section~\ref{sec:cycle}).
\end{enumerate}
The permanent lemmas share one idea---divisibility of a permanent by the order of a
free group acting on the index set while preserving the summand---which we isolate
as~(5) and reuse throughout.

\section{Honest assessment}\label{sec:assessment}

We state the limitations plainly, because the paper's credibility depends on it.
Individually the four discoveries are minor (on a grand-scheme scale we would rate
them $2$--$4$ out of $10$): they concern niche objects, and where they touch deep
structure---Case~2, the Ramanujan-coverings framework---that
structure was established by others in 2018, capping the novelty. Rigor and
machine-checking do \emph{not} buy significance; a formally verified statement
about a shallow object is still shallow. Case~3's central number-theory
question is open, and our own analysis suggests its natural proof targets are
equivalent to it rather than easier.

The claim is therefore methodological. What the pipeline delivers that a
pen-and-paper note does not is: (a) a reproducible, kernel-checked realization of
the compute-conjecture-prove loop, end to end, on discovered (not imported)
objects; (b) the first formally verified instance of the Ramanujan-coverings circle
of ideas; (c) a reusable formal permanent library Mathlib lacked; and (d) an
enforced discipline of retraction when a computation's proxy proves unfaithful.
Not a notebook of new deep mathematics---a proof of concept that Ramanujan's method
scales to \emph{rigor}, if not yet to \emph{depth}.

\begin{thebibliography}{9}
\bibitem{HPS} C.~Hall, D.~Puder, W.~F.~Sawin,
  \emph{Ramanujan coverings of graphs}, Adv.\ Math.\ \textbf{323} (2018),
  367--410.
\bibitem{MSS} A.~Marcus, D.~Spielman, N.~Srivastava,
  \emph{Interlacing families~I: bipartite Ramanujan graphs of all degrees},
  Ann.\ of Math.\ \textbf{182} (2015), 307--325.
\bibitem{HL} O.~J.~Heilmann, E.~H.~Lieb,
  \emph{Theory of monomer-dimer systems}, Comm.\ Math.\ Phys.\ \textbf{25}
  (1972), 190--232.
\bibitem{Smith} H.~J.~S.~Smith,
  \emph{On the value of a certain arithmetical determinant},
  Proc.\ London Math.\ Soc.\ \textbf{7} (1875--76), 208--212.
\bibitem{Ryser} H.~J.~Ryser, \emph{Combinatorial Mathematics},
  Carus Math.\ Monographs \textbf{14}, MAA, 1963.
\bibitem{Mathlib} The mathlib Community,
  \emph{The Lean mathematical library}, CPP 2020, 367--381.
\end{thebibliography}

\end{document}
